\section{Supported asymptotic complexes}
\label{sec:asymptotic}

The goal is an analogue of Lemma~\ref{ds:ordinary-norm} for sequences of
bounded free complexes modulo free sequences of negligible normalized
rank, proved in Theorem~\ref{as:norm}. Two points require care: support
may arise only after passing to the quotient, and supported objects may
be idempotent summands whose projectors do not lift before quotienting.
Their homology lengths therefore cannot in general be computed from the
raw homologies of arbitrary lifts. We first construct nonnegative lengths
on uniformly supported lifts, extend them through the supported quotient
and its idempotent completion, and prove their compatibility with scalar
extension on every supported summand.

Throughout this section, $(A,J)$ is an excellent regular local
$\Z_{(p)}$-algebra of dimension $e>0$ with residue characteristic $p$,
where $J$ is the maximal ideal, and $z_1,\ldots,z_e$ is a regular system
of parameters. Let $T$ be a finite
local domain over $A$ with the same residue field, and put
$r=\rank_A T$. Let $C$ be a $T$-algebra equipped with a nonnegative
additive length $\lambda$ on its modules annihilated by powers of $J$.
The assumptions on this length that are used in this section are
\begin{equation}
 \lambda H_j(z_1^v,\ldots,z_e^v;C)=0\quad(j>0),
 \qquad
 \lambda\bigl(C/(z_1^v,\ldots,z_e^v)C\bigr)=v^e r
 \quad(v\geq 1).
 \label{as:length-hypotheses}
\end{equation}
The length is allowed to vanish on nonzero modules. In the applications,
Theorem~\ref{len:input} supplies these hypotheses, as well as the stronger
parameter properties used in the other sections. Write $K_v$ for the
Koszul complex $K(z_1^v,\ldots,z_e^v;A)$, in homological degrees
$0,\ldots,e$.

\subsection{The rank quotient and its support condition}

Fix positive integers $s_n$ and a nonprincipal ultrafilter $\mathcal U$
on the positive integers. Rank estimates $O(s_n)$ and $o(s_n)$ always
refer to the ordinary limit, not to $\mathcal U$. If $(a_n/s_n)$ is
bounded, write
\[
 \overline a=\lim_{n\to\mathcal U}\frac{a_n}{s_n}.
\]
For $S=A$ or $T$, consider the additive category of sequences
$(S^{a_n})_n$ with $a_n=O(s_n)$, whose morphisms are arbitrary
sequences of matrices over $S$. Let $\mathcal D_S$ be its perfect
hull: take complexes in one common bounded degree range for the whole
sequence, and then idempotent completion. Let
$\mathcal N_S$ be the thick subcategory generated by the free
sequences with $a_n=o(s_n)$, and define
\[
 \mathcal P_S=(\mathcal D_S/\mathcal N_S)^{\natural}.
\]
Here $(-)^{\natural}$ denotes idempotent completion. The quotient
functor is denoted by $q_S$.

A thick subcategory is closed under shifts, cones, and direct summands.
The Verdier quotient inverts maps whose cones lie in that subcategory;
idempotent completion then adjoins the images of idempotent endomorphisms.

An object $X$ in $\mathcal D_S$ or $\mathcal P_S$ is called
\emph{supported on $J$} if there is an integer $v\geq1$ such that
$z_j^v\operatorname{id}_X=0$ for every $j$. The identity is taken
in the category containing $X$. Thus in $\mathcal D_S$ this is a
uniform homotopical support condition on the sequence, whereas in
$\mathcal P_S$ the actions need vanish only after the rank-small
quotient. Write $\mathcal S_S\subseteq\mathcal D_S$ and
$\mathcal P_S^J\subseteq\mathcal P_S$ for the resulting full
subcategories. They are thick: nilpotence of central actions is
preserved by shifts, cones, and retracts, allowing larger powers
for cones. We identify this condition with central-localization
support before applying descent at the end of the section.

\begin{remark}
\label{as:nonuniform-warning}
Raw homology lengths need not survive the rank quotient. Over a
discrete valuation ring $A=T$ with parameter $z$, take $s_n=n$ and
the rank-one free resolution of $T/(z^n)$. This sequence belongs to
$\mathcal N_T$, whereas its raw normalized degree-zero homology
length is $1$. It has no uniform annihilating power before
quotienting. For each fixed $v$, however, its homology lengths after
tensoring with $K(z^v;T)$ are bounded independently of $n$ and hence
become zero after normalization. This is why we construct homology
lengths on uniformly supported lifts first, and then extend them to
supported quotient objects.
\end{remark}

\subsection{The supported subquotient}

The next lemma supplies that extension's categorical starting point:
every supported quotient object is a summand of a uniformly supported
lift, and morphisms between such lifts can be computed using supported
roofs. Neither assertion requires the summand projector to lift.

\begin{lemma}
\label{as:supported-quotient}
The natural exact functor
\[
 \mathcal S_S/(\mathcal S_S\cap\mathcal N_S)
 \longrightarrow \mathcal D_S/\mathcal N_S
\]
is fully faithful. Its idempotent completion is
$\mathcal P_S^J$.
\end{lemma}

\begin{proof}
First let $X\in\mathcal S_S$, $Y\in\mathcal N_S$, and
$f:X\to Y$. Choose $v$ such that $z_j^v\operatorname{id}_X=0$
for every $j$. Centrality gives $z_j^vf=0$. The exact Hom sequence
of the triangle
\[
 \operatorname{fib}(z_j^v:Y\to Y)\longrightarrow Y
 \xrightarrow{z_j^v}Y
\]
therefore lifts $f$ through the first object. Repeat this operation
for all the parameters, always using the zero action on the original
source $X$. The result factors $f$ through
$K_v\otimes_A Y[-e]$. This object belongs to $\mathcal N_S$ by
thickness, and to $\mathcal S_S$ by the usual Koszul homotopies.
Thus every map from a supported object to a small object factors
through a supported small object.

Here is the resulting roof argument. A morphism between $X,Y$ in
$\mathcal S_S$, computed in $\mathcal D_S/\mathcal N_S$, has a
representative
\[
 X\xleftarrow{s}Z\xrightarrow{f}Y,
 \qquad \operatorname{cofib}(s)=N_0\in\mathcal N_S.
\]
In the triangle $Z\to X\xrightarrow{u}N_0$, factor $u$ through
$N_1\in\mathcal S_S\cap\mathcal N_S$ by the preceding paragraph.
Let $Z_1$ be the fiber of $X\to N_1$. The induced morphism of
triangles gives $c:Z_1\to Z$ and $s_1:Z_1\to X$ with $sc=s_1$.
Now $Z_1\in\mathcal S_S$ and $\operatorname{cofib}(s_1)=N_1$,
so $X\xleftarrow{s_1}Z_1\xrightarrow{fc}Y$ is a representative
entirely in the supported subquotient. This proves fullness.

For faithfulness, suppose a map $f:Z\to Y$ between supported
objects becomes zero in $\mathcal D_S/\mathcal N_S$. The calculus
of Verdier roofs gives a map $t:W\to Z$ with small cone and $ft=0$.
The triangle of $t$ makes $f$ factor through that small cone.
Factoring its first leg through a supported small object as above
shows that $f$ is already zero in the supported subquotient. Apply
this to the numerator of any roof to obtain faithfulness.

Finally take $X\in\mathcal P_S^J$. It is a retract of $q_S(Y)$
for some $Y\in\mathcal D_S$. Choose $v$ with all $z_j^v$ acting
as zero on $X$. The cone of each such zero map splits, so $X$ is a
retract of $K_v\otimes_A X$. The latter is a retract of
$q_S(K_v\otimes_A Y)$, whose lift is uniformly supported. Hence
$X$ belongs to the idempotent completion of the fully faithful
supported subquotient. The reverse inclusion follows because
support is closed under retracts.
\end{proof}

\subsection{Nonnegative homology lengths}

Every supported quotient object is now accessible through a uniformly
supported lift and a summand. We next construct its nonnegative homology
lengths and show that they give additive Euler evaluations.

For ordinary length, let $\mathcal A_T$ be the abelian category
of sequences of finite $T$-modules annihilated by some common power
of $J$ and with lengths $O(s_n)$. For the $C$ evaluation, let
$\mathcal A_C$ be the category of sequences of $C$-modules with
the same uniform torsion condition and with finite $\lambda$-lengths
$O(s_n)$. The power and the implied constant can depend on the
object. Kernels, cokernels, and extensions preserve these conditions,
so both categories are abelian. In each case the sequences of
normalized ultralimit length zero form a Serre subcategory. Denote
the Serre quotients by $\overline{\mathcal A}_T$ and
$\overline{\mathcal A}_C$, and their nonnegative additive lengths
by $\overline\length_T$ and $\overline\lambda$, respectively.
These lengths are finite on every object.
Here negligible length means zero ultralimit; the free sequences
generating $\mathcal N_T$ still satisfy the stronger ordinary estimate
$o(s_n)$. We will show that supported objects killed by the rank
quotient have zero homology in these abelian quotients.

\begin{proposition}
\label{as:evaluations}
There are homological functors
\[
 \mathsf H_i^{\mathrm{ord}}:\mathcal P_T^J
       \longrightarrow\overline{\mathcal A}_T,
 \qquad
 \mathsf H_i^C:\mathcal P_T^J
       \longrightarrow\overline{\mathcal A}_C
\]
whose values on a uniformly supported lift are its ordinary homology
sequence and its homology sequence after tensoring with $C$,
respectively. For each object $X$, their lengths
\[
 b_i^{\mathrm{ord}}(X)=\overline\length_T
       \mathsf H_i^{\mathrm{ord}}(X),
 \qquad
 b_i^C(X)=\overline\lambda\mathsf H_i^C(X)
\]
are finite nonnegative real numbers, zero outside a finite range.
Their Euler sums define homomorphisms
\[
 \chi_{\mathrm{ord}},\chi_C:K_0(\mathcal P_T^J)\longrightarrow\mathbb R.
\]
If $F=(F_n)$ is a bounded free sequence whose image is supported,
and $v$ makes all $z_j^v$ act as zero on that image, then for either
choice of evaluation, writing $b_i(F)=b_i(q_TF)$,
\begin{equation}
 (1+t)^e\sum_i b_i(F)t^i
  =\sum_i t^i\lim_{n\to\mathcal U}
     \frac{\operatorname{len}
       H_i(K_v\otimes_A F_n)}{s_n}.
 \label{as:koszul-polynomial}
\end{equation}
Here $\operatorname{len}=\length_T$ in the ordinary case; in the
$C$ case the complex on the right is further tensored over $T$
with $C$, and $\operatorname{len}=\lambda$. The identity is one
of finite Laurent polynomials.
\end{proposition}

\begin{proof}
We first bound homology on $\mathcal S_T$, then show that it kills
$\mathcal S_T\cap\mathcal N_T$, and finally extend it to the
idempotent-completed quotient.

Begin with a bounded free sequence $F$ in $\mathcal S_T$. Choose
uniform powers annihilating its identity. Splitting their cones
gives an isomorphism in the homotopy category
\begin{equation}
 K_v\otimes_A F\simeq
 \bigoplus_{j=0}^e F[j]^{\oplus\binom ej}.
 \label{as:split-koszul}
\end{equation}
For ordinary length, every Koszul homology $H_j(K_v\otimes_A T)$
has finite length. Taking Koszul homology first in the bounded
double complex bounds the homology lengths of $K_v\otimes_A F_n$
by fixed constants times the ranks of the terms of $F_n$. They
are therefore $O(s_n)$. For $C$, the same spectral sequence uses
Equation~\eqref{as:length-hypotheses}: positive Koszul rows have
length zero and the zero row has finite length. Additivity and
nonnegativity give the same bound. The splitting in
Equation~\eqref{as:split-koszul} bounds the homologies of $F_n$
themselves. They have a common primary annihilator, because a power
of every $z_j$ annihilates them.

A general object $X\in\mathcal S_T$ may be a retract in
$\mathcal D_T$ of a bounded free sequence $F$ that is not supported.
Choose $v$ killing the parameter actions on $X$. Then $X$ is a retract
of $K_v\otimes_A X$, which is a retract of the uniformly supported
bounded free sequence $K_v\otimes_A F$. The preceding bounds apply
to this sequence and hence to the homology summands defining $X$.
This is a retract before taking the rank quotient. Thus homology on
$\mathcal S_T$ takes
values in the indicated abelian length quotients.

Fix $v$. Koszul tensoring a rank-small free generator has homology
of normalized length zero by the same bounds. The class of objects
$Y\in\mathcal D_T$ with this property for $K_v\otimes_A Y$
is thick: use homology long exact sequences, nonnegativity, and
retracts. It therefore contains $\mathcal N_T$. If
$X\in\mathcal S_T\cap\mathcal N_T$, choose $v$ killing its
parameter actions. Then $X$ is a retract of $K_v\otimes_A X$,
so its homology is zero in both length quotients.

The homology functors consequently invert maps whose cones are
supported and small. They descend to
$\mathcal S_T/(\mathcal S_T\cap\mathcal N_T)$, preserving
homology long exact sequences. Now consider summands after quotienting.
By Lemma~\ref{as:supported-quotient},
every object of $\mathcal P_T^J$ is a summand of an object of this
subcategory. Extend a homology functor to a summand $(X,a)$ by
taking $\operatorname{im}\mathsf H_i(a)$ in its abelian target.
This does not require a sequence of chain projectors lifting $a$.
Distinguished triangles in an idempotent completion are summands
of distinguished triangles in the original triangulated category
\cite[Definition~1.10 and Theorem~1.12]{BalmerSchlichting}.
Taking images of compatible idempotents preserves exact sequences
in an abelian category, so the extended functors are still
homological.

Each object is a summand of a fixed bounded uniformly supported
Koszul lift. Its homology objects therefore vanish outside a finite
range, and their lengths are bounded by those of that lift. This
proves finiteness, nonnegativity, and boundedness of the $b_i$.
The finite long exact sequences prove additivity of the Euler sums.

For the last assertion, the splitting in
Equation~\eqref{as:split-koszul} now takes place in $\mathcal P_T$,
since the chosen powers act as zero there. Applying the extended
homological functors gives the left side of
Equation~\eqref{as:koszul-polynomial}. Its Koszul tensor has a
uniformly supported lift, whose actual homology computes the right
side. This proves the identity without defining its coefficients
by formal division.
\end{proof}

In particular, Equation~\eqref{as:koszul-polynomial} assigns zero
to the quotient object in Remark~\ref{as:nonuniform-warning}, as
required. Its nonnegative coefficients come from homological
functors, not from an assumed positivity of a polynomial quotient.

\subsection{Scalar extension and nonlifting idempotents}

Construct $\overline{\mathcal A}_A$ and its length in the same way
using ordinary finite $A$-modules. The preceding construction over
$A$ with ordinary length defines $\chi_A$ on $\mathcal P_A^J$.
Scalar extension gives an exact functor
$\mathcal P_A^J\to\mathcal P_T^J$.

The supported Euler evaluations are defined on every object of
$\mathcal P_T^J$. To apply descent, it remains to prove that they agree
on all classes coming from $\mathcal P_A^J$, including idempotent
summands.

The normalized evaluation uses the following consequence of the
length hypotheses. We state it for the present base $A$, which need
not be complete.

\begin{lemma}
\label{as:finite-tor}
For every finite-length $A$-module $V$,
\[
 \lambda\Tor_j^A(V,C)=0\quad(j>0),
 \qquad
 \lambda(V\otimes_A C)=r\length_A V.
\]
All these modules are annihilated by a power of $J$.
\end{lemma}

\begin{proof}
The Koszul and composition-series proof of
Lemma~\ref{ds:finite-length-basechange} applies here: it uses regularity
and the length hypothesis, but not completeness. A power of $J$
annihilating $V$ also annihilates its derived tensor homology.
\end{proof}

\begin{proposition}
\label{as:base-equality}
For every $X\in\mathcal P_A^J$,
\[
 \chi_{\mathrm{ord}}(X\otimes_A^{\mathbf L}T)
   =r\chi_A(X)
   =\chi_C(X\otimes_A^{\mathbf L}T).
\]
These equalities hold for summands defined by idempotents in the
Verdier quotient, without assuming that the idempotents lift to
chain maps before localization.
\end{proposition}

\begin{proof}
The two evaluations recover the factor $r$ differently: the ordinary
one uses the alternating rank of a resolution of $T$, while the
normalized one uses the vanishing of positive Tor lengths. To make
both calculations on an arbitrary summand, we descend their finite
spectral sequences through supported Verdier roofs before taking
idempotent images.

Let $Y$ be a uniformly supported bounded free sequence over $A$ and
let $a$ be an idempotent endomorphism of its image in
$\mathcal S_A/(\mathcal S_A\cap\mathcal N_A)$. Every supported
object of $\mathcal P_A$ has this form as a summand, by
Lemma~\ref{as:supported-quotient}: in its proof one may enlarge the
initial lift to a bounded free sequence before taking the Koszul
tensor. Put
\[
 V_b=\operatorname{im}\mathsf H_b^A(a)
       \quad\text{in }\overline{\mathcal A}_A.
\]
These images exist in the abelian quotient even when $a$ has no
chain-level projector representing it.

Choose a fixed bounded finite free $A$-resolution $P_\bullet$ of
$T$; such a resolution exists because $A$ is regular. The double
complex $Y\otimes_A P_\bullet$ computes
$Y\otimes_A^{\mathbf L}T$. Filtering by the degree of $P$ gives
a finite spectral sequence with first page
\begin{equation}
 E^1_{u,b}=H_b(Y)\otimes_A P_u.
 \label{as:ordinary-page}
\end{equation}
It is natural in homotopy-category maps from this page onward:
homotopies in $Y$ preserve the filtration and act trivially on
the first page. On retracts in the prequotient category the same
construction is defined by taking the corresponding idempotent
images on every page and on the abutment filtration. Its terms
and finite abutment filtration can be
regarded as $A$-modules. The abutment's ordinary $A$-length equals
its $T$-length, since the residue fields agree.

A map with supported small cone is an isomorphism on $H_b(Y)$
in $\overline{\mathcal A}_A$. It is thus an isomorphism on
Equation~\eqref{as:ordinary-page}, on every subsequent page, and
on the finite filtered abutment in the length quotient. The
spectral sequence therefore descends through all Verdier roofs
in the supported subquotient. The natural augmentation
$\operatorname{Tot}(Y\otimes_A P_\bullet)\to Y\otimes_A T$
identifies its abutment with the homology of scalar extension,
and this identification also descends through those roofs.
Thus the induced action of $a$ on the abutment is precisely the
scalar-extended idempotent, not a separately chosen action.

In particular $a$ acts as an actual
idempotent on every page and on that filtered abutment. The
differentials commute with this action. Taking its images commutes
with passage to homology, since an idempotent decomposes a complex
in an abelian category into its image and complementary image.
It also preserves the finite exact sequences defining the
abutment filtration: for a filtration preserved by an idempotent
$a$, one has $\operatorname{im}(a)\cap F_u=a(F_u)$ and hence
$\operatorname{gr}_u\operatorname{im}(a)
=\operatorname{im}(\operatorname{gr}_u a)$.
Consequently the Euler length on the
scalar-extended summand is the Euler length of the images on
the first page. They are $\rank_A(P_u)$ copies of $V_b$, and hence
this value is
\begin{equation}
 \left(\sum_u(-1)^u\rank_A P_u\right)
 \sum_b(-1)^b\overline\length_A V_b
 =r\sum_b(-1)^b\overline\length_A V_b.
 \label{as:ordinary-idempotent}
\end{equation}
The equality of alternating ranks with $r$ follows by tensoring
the resolution with the fraction field of $A$.

For $C$, use the finite hyper-Tor spectral sequence
\begin{equation}
 {}^CE^2_{u,b}=\Tor_u^A(H_b(Y),C)
 \ \Longrightarrow\ H_{u+b}(Y\otimes_A^{\mathbf L}C).
 \label{as:c-page}
\end{equation}
It can be constructed from the functorial finite Postnikov
filtration of $Y$. It has bounded $b$ and $0\leq u\leq e$, so
there is no unbounded convergence issue. Lemma~\ref{as:finite-tor}
makes every positive Tor term zero in
$\overline{\mathcal A}_C$. It also implies that
$V\mapsto V\otimes_A C$ induces an exact functor
\[
 \overline{\mathcal A}_A\longrightarrow
 \overline{\mathcal A}_C
\]
that multiplies length by $r$. Indeed, on a short exact sequence
its possible failure of left exactness is the image of a
positive Tor module of zero length. Negligible-length sequences
map to negligible-length sequences by the same lemma. This
explains both exactness after the Serre quotient and descent of
the functor to that quotient.

It follows that a supported-small-cone map induces isomorphisms
on the second page of Equation~\eqref{as:c-page} and its finite
filtered abutment in the length quotient. Thus this spectral
sequence too is functorial through the supported Verdier roofs.
Taking the images of the same idempotent $a$ gives only the zero
Tor row, namely the images of $V_b$ under this exact tensor
functor. Their Euler length is
$r\sum_b(-1)^b\overline\length_A V_b$, exactly the value in
Equation~\eqref{as:ordinary-idempotent}. Finally this last sum
is $\chi_A(Y,a)$ by construction. This proves both equalities
for every summand and hence for every $X\in\mathcal P_A^J$.
\end{proof}

\subsection{Connective coefficients and the norm comparison}

We have constructed the supported Euler evaluations and proved their
agreement after scalar extension from $A$. To extend that agreement
to conjugate sums, we now put the rank quotient in the form required
by Theorem~\ref{ds:descent}. The coefficient algebra must be
connective, and the support condition must be central-localization
support.

\begin{lemma}[Connective coefficients]
\label{as:connective}
There is a connective associative derived $A$-algebra $H$ and
compatible equivalences
\[
 \mathcal P_A\simeq\Perf(H),
 \qquad
 \mathcal P_T\simeq\Perf(H\otimes_A^{\mathbf L}T).
\]
The equivalences respect central localization by elements of $A$.
\end{lemma}

\begin{proof}
The sequence $E_S=(S^{s_n})_n$ is an additive generator. Indeed, if
$a_n=O(s_n)$, there is a single integer $c$ such that $a_n\leq cs_n$
for every $n$, after increasing $c$ to handle the finitely many
exceptional indices. Coordinate inclusions and projections make
$(S^{a_n})_n$ a summand of $E_S^c$. Consequently
\[
 \mathcal D_S\simeq\Perf(B_S),
 \qquad B_S=\prod_n M_{s_n}(S).
\]

The $A$-module underlying $B_A$ is flat. One can check this by the
ideal criterion: every finitely generated ideal of the Noetherian
ring $A$ is finitely presented, so its tensor product with this
product of finite free modules is the corresponding product of its
tensor products; each required injection is then a product of
injections. Moreover, $T$ is finitely presented over $A$, so tensoring
with $T$ commutes with that product. Thus
\begin{equation}
 B_A\otimes_A^{\mathbf L}T
 \simeq B_A\otimes_A T
 \cong B_T.
 \label{as:product-base-change}
\end{equation}

The same product argument makes every morphism module between the
original free sequences flat over $A$. Hence the explicit dg quotient
of this degree-zero category by its full rank-small subcategory computes
the derived quotient \cite[Section~3.1 and Theorem~3.4]{Drinfeld}.
Its morphisms are sums of paths formed from tensor products of the
original morphism modules and contracting homotopies in homological
degree $1$. They are concentrated in nonnegative degrees and degreewise
$A$-flat, hence are $K$-flat. Passing to the perfect hull and idempotent
completion kills exactly $\mathcal N_A$. The image of $E_A$ remains a
generator, and its endomorphism algebra
$H=\operatorname{End}_{\mathcal P_A}(q_AE_A)$ is connective, giving
$\mathcal P_A\simeq\Perf(H)$. This assertion concerns the chosen
generator, not arbitrary shifts of perfect objects.

Since $T$ is finitely presented over $A$, tensoring each original
morphism module with $T$ gives the corresponding module of morphisms
between free $T$-sequences. Tensoring the path formula therefore gives
exactly the $T$-linear dg quotient by the rank-small free sequences.
The same flatness criterion applies over the Noetherian ring $T$.
The upstairs endomorphism algebra is consequently
$H\otimes_A T\simeq H\otimes_A^{\mathbf L}T$, and its perfect hull
and idempotent completion are $\mathcal P_T$. This also agrees with
the preservation of exact quotients under scalar extension
\cite[Lemma~3.12]{LT}. All constructions are centrally $A$-linear and
preserve central localization.
\end{proof}

Under these equivalences, the support condition used throughout this
section is exactly the one in Theorem~\ref{ds:descent}. Indeed, for
a perfect module $X$, localization at a central element $z_j$ is
the telescope of multiplication by $z_j$ in the module category.
Compactness implies that $X[1/z_j]=0$ if and only if some power
$z_j^v$ kills its identity. A common power can be chosen for the
finitely many parameters. The same argument for $\Perf(B_S)$
identifies $\mathcal S_S$ with the supported objects before the
rank quotient. In particular, this comparison does not require a
uniformly supported lift of an object in $\mathcal P_S^J$.

Suppose now that a finite group $G$ acts on $T$ over $A$, that
$\Spec T\to\Spec A$ is surjective, and that $G$ is transitive
on every geometric fiber, as in Theorem~\ref{ds:descent}. It acts
on $\mathcal P_T^J$ by scalar twisting. We write $gX$ for the
functor induced by $g^*P=P\otimes_{T,g}T$, with the convention of
Section~\ref{sec:descent}.

\begin{theorem}[Asymptotic norm comparison]
\label{as:norm}
Let $A,J,T,C,\lambda$ satisfy the setup of this section and
\eqref{as:length-hypotheses}. Suppose a finite group $G$ acts on $T$
over $A$, the map $\Spec T\to\Spec A$ is surjective, and $G$ is
transitive on every geometric fiber. Then, for every
$X\in\mathcal P_T^J$,
\[
 \sum_{g\in G}\chi_{\mathrm{ord}}(gX)
   =\sum_{g\in G}\chi_C(gX).
\]
Neither invariance of $\lambda$ under $G$ nor equality of the two
evaluations on individual classes is assumed.
\end{theorem}

\begin{proof}
Lemma~\ref{as:connective} identifies the two supported categories
with the supported perfect categories of $H$ and
$H\otimes_A^{\mathbf L}T$, where $H$ is connective.
The identification over $T$ is compatible with the $G$-action:
scalar twisting acts on the $T$ factor in the coefficient algebra.
Theorem~\ref{ds:descent} therefore gives a surjection
\[
 K_0(\mathcal P_A^J)\otimes\Q
 \longrightarrow
 \bigl(K_0(\mathcal P_T^J)\otimes\Q\bigr)^G.
\]
The norm class $\sum_g[gX]$ is invariant. The homomorphism
$\chi_{\mathrm{ord}}-\chi_C$, extended rationally to the
Grothendieck group, vanishes on the image from $A$ by
Proposition~\ref{as:base-equality}. Applying it to a preimage
of this norm class proves the assertion.
\end{proof}
